% ID: SEL-D0050-OB-E04
% ESTADO: cotejada. Segunda pasada visual y borrador compilado revisados.
% Fuente: Selectividad/sel_2013_matematicas_aplicadas/junio/titular/EXAMEN 6 -ANDALUCÍA.pdf
% SHA-256: 1f0983e351b8435c5b6c7b68b05fe2a64034faf77a85a1f9480407aa0f2a88aa
% Página: 2; opción B; ejercicio original 4.

El tiempo que los españoles dedican a ver la televisión los domingos es una variable aleatoria que sigue una distribución Normal de media desconocida y desviación típica 75 minutos. Elegida una muestra aleatoria de españoles se ha obtenido, para la media de esa distribución, el intervalo de confianza $(188.18,208.82)$, con un nivel del 99\%.
\begin{enumerate}[label=\alph*)]
\item \textbf{(1.5 puntos)} Calcule la media muestral y el tamaño de la muestra.
\item \textbf{(1 punto)} Calcule el error máximo permitido si se hubiese utilizado una muestra de tamaño 500 y un nivel de confianza del 96\%.
\end{enumerate}
